\section{Suffering, Secondary Causes, and What May Be Said}

The doctrine has now been stated and fenced. What remains is the case it exists for, and the case is a person in pain. Everything in the preceding sections bears on how one may and may not speak to that person---which is not a change of subject, because a doctrine of providence that cannot be spoken truthfully at a bedside has failed at the only point where it was ever going to be tested.

\subsection{What Providence Does Not License Anyone to Say}

The Catechism's paragraph on the scandal of evil begins by refusing the short answer. Why does evil exist, if God cares for all his creatures? ``To this question, as pressing as it is unavoidable and as painful as it is mysterious, no quick answer will suffice.'' What follows is not a theodicy but a list of the whole Christian mystery---creation's goodness, the drama of sin, the covenants, the Incarnation, the Spirit, the Church, the sacraments, the call to blessedness---with the comment that ``there is not a single aspect of the Christian message that is not in part an answer to the question of evil.''\endnote{CCC 309, official English checked 2026-07-25. The paragraph's refusal of a quick answer is doctrinally significant: it means that no formula, including any in this article, is offered as the resolution of the difficulty.}

Two further paragraphs draw the limits inside which consolation may operate. God ``can bring a good from the consequences of an evil, even a moral evil, caused by his creatures''---the text's example is Joseph's ``you meant evil against me; but God meant it for good''---and from the greatest moral evil ever committed, the murder of the Son, God brought the greatest of goods. Then the sentence that stops the argument from running away: ``But for all that, evil never becomes a good.'' And finally, the epistemic limit: the ways of providence ``are often unknown to us,'' and only at the end will we know the ways by which God has guided his creation.\endnote{CCC 312 and 314, official English checked 2026-07-25.}

Put together, these yield a rule of speech that is stricter than most pious conversation. One may say that God permitted this and did not cause the fault in it; that he can and will draw good from it; that the good is promised and not yet displayed; and that the sufferer is not being told what the good is, because nobody knows. One may not say that this happened \emph{in order to} teach the sufferer a particular lesson, or that the death of a child was the means to some identified benefit, or that the affliction is a punishment. The last is not merely tactless: the Gospel rejects the inference in terms when the disciples try it, and Aquinas's careful distinction---that natural corruption or punishment may be willed incidentally under a good---is a statement about the order of the universe, not a diagnostic tool for individual biographies.

\subsection{The Joseph Pattern}

Genesis 50:20 is the tradition's paradigm, and its structure is worth extracting because it is the structure of every honest use of the doctrine. Three things are asserted in one sentence and none is retracted. The brothers meant evil: their intention was wicked, and Joseph names it as such thirty years later, to their faces. God meant it for good: the divine intention ran through the same event without endorsing theirs. And the good is specified only in retrospect and only in part---that many people should be kept alive.

Notice what Joseph does not do. He does not tell his brothers that their crime was fine because it worked out; he tells them not to be afraid, which is forgiveness, not exculpation. He does not claim to have known during his years in the pit and the prison what the outcome would be. And his statement of the good is public and verifiable---a famine survived---not a private spiritual benefit assigned to him by someone else.

The pattern therefore licenses retrospective and partial reading of one's own history, in the light of an outcome one can actually see. It does not license prospective assignment of meaning to someone else's suffering. The difference between those two operations is most of the difference between the consolation of the saints and the cruelty of the comforters of Job.

\subsection{Lament Is Not a Failure of Surrender}

A doctrine of surrender that forbids complaint would have to excise the Psalter, most of Job, Jeremias, and the cry from the cross. It does not forbid it. Job's ``if we have received good things at the hand of God, why should we not receive evil?'' is spoken by a man who will spend thirty chapters demanding an answer from God and who is vindicated at the end against the friends who defended providence with better manners and worse theology.

Nor does the tradition require that surrender be felt. De Sales, describing acquiescence in a death he could not prevent, has his speaker consent ``in the point of my spirit, in spite of all the opposition of the inferior powers of my soul.''\endnote{\emph{Treatise}, Book~IX, ch.~VI, in Mackey's translation, as cited above. The phrase locates consent in the will's summit and expressly allows the rest of the person to be in revolt---a distinction that protects the doctrine from being falsified by grief, and protects the griever from being told that grief is infidelity.} That is a precise anthropological claim, and it does considerable pastoral work: consent is an act of the will, not a state of the emotions, and a person weeping is not thereby a person refusing.

\subsection{Why Secondary Causes Matter Pastorally}

The architecture set out earlier is not decoration on the consolation; it is what makes the consolation sayable. If God were the only cause, then every tumour would be a direct divine act, every assault a divine gesture, and the sufferer's anger would have exactly one legitimate object. Because created causes are real---because fire really heats, cells really divide, and men really choose---the world's inventory of harms is not a list of divine decisions. It is the behaviour of a real, wounded, contingent order that God holds in being and governs to an end.

That is also why the language of ``God's plan,'' used carelessly, does more damage than the doctrine it garbles. Providence is not a script in which each line has been written for its dramatic effect. It is a plan of ordering: an intellect that knows every detail, holds every cause in being, permits what it does not will, and guarantees an end that is not visible from here. What follows for the sufferer is not that the tumour was chosen for him, but that it did not escape, that it is not the last word, and that he is not alone inside it.
